\documentclass[11pt]{article}
\usepackage{rabbithole}

\phasenumber{11}
\phasetitle{Configuration and automation}
\phasesubtitle{Aliases, attributes, and hooks that stop a bad commit before it exists.}

\hypersetup{pdftitle={Rabbit Holes, Phase 11: Configuration and automation},
            pdfauthor={Accelerate, Manipal University Jaipur},
            pdfsubject={git and GitHub handbook}}

\begin{document}
\maketitlepage

\section{Four levels of config}

Git reads settings from several files and the most specific wins.

\vspace{0.5em}
\begin{tabularx}{\linewidth}{@{}L{2.6cm} L{4.4cm} X@{}}
\toprule
\rhtablehead{Level} & \rhtablehead{File} & \rhtablehead{Applies to} \\
\midrule
\co{-{}-system} & \co{/etc/gitconfig} & Every user on the machine \\
\co{-{}-global} & \co{\textasciitilde/.gitconfig} & You, everywhere \\
\co{-{}-local} & \co{.git/config} & One repository \\
\co{-{}-worktree} & \co{.git/config.worktree} & One worktree \\
\bottomrule
\end{tabularx}
\vspace{0.5em}

\begin{shell}
git config --list --show-origin
git config --get user.email
git config --global --edit
\end{shell}

\co{--show-origin} is the one to remember. When a setting is not doing what you
expect, it tells you which file it came from, and the answer is usually a
\co{--local} value you set months ago and forgot.

\subsection{Different identities for different work}

Personal projects under one email, university work under another:

\begin{terminal}
[includeIf "gitdir:~/work/"]
    path = ~/.gitconfig-work
\end{terminal}

Put that in \co{~/.gitconfig}, put the work identity in \co{~/.gitconfig-work}, and
git switches automatically based on where the repository lives. The trailing slash
matters.

\section{Settings worth turning on}

\begin{shell}
git config --global init.defaultBranch main
git config --global pull.rebase true
git config --global fetch.prune true
git config --global rebase.autoSquash true
git config --global rerere.enabled true
git config --global merge.conflictStyle zdiff3
git config --global diff.algorithm histogram
git config --global push.autoSetupRemote true
git config --global column.ui auto
git config --global help.autocorrect 20
\end{shell}

Three of those deserve a note.

\co{push.autoSetupRemote} means \co{git push} works on a brand new branch without
\co{-u origin <name>}. It removes one of the most common small frictions in daily
use.

\co{diff.algorithm histogram} produces noticeably better diffs on code that has been
moved or re-indented, at no real cost.

\co{help.autocorrect 20} makes git run the obvious correction after two seconds when
you typo a command. Set it to \co{0} if you find that alarming rather than helpful.

\subsection{Line endings}

The classic cross-platform problem: Windows uses CRLF, everything else uses LF, and
without configuration every file shows as entirely modified when the other person
touches it.

\begin{shell}
git config --global core.autocrlf true
git config --global core.autocrlf input
\end{shell}

\co{true} on Windows, \co{input} on macOS and Linux.

\begin{note}
Better still, do not rely on each person configuring it. Put a \co{.gitattributes}
file in the repository with \co{* text=auto eol=lf} and the rule travels with the
project, applying to everyone regardless of their settings. This repository does
exactly that.
\end{note}

\section{Aliases}

\begin{shell}
git config --global alias.st status
git config --global alias.co checkout
git config --global alias.br branch
git config --global alias.lg "log --graph --oneline --decorate --all"
git config --global alias.last "log -1 HEAD --stat"
git config --global alias.unstage "restore --staged"
git config --global alias.amend "commit --amend --no-edit"
\end{shell}

Aliases starting with \co{!} run a shell command instead, from the repository root:

\begin{shell}
git config --global alias.root "!pwd"
git config --global alias.cleanup "!git branch --merged main | grep -v main | xargs -r git branch -d"
git config --global alias.sync "!git fetch --prune && git pull --rebase"
\end{shell}

That \co{cleanup} one deletes every local branch already merged into \co{main},
which after an active month is most of them. Read it before you run it, which is
good advice for any alias you copy from a document.

\begin{tryit}
Add this one and use it for a week:

\begin{shell}
git config --global alias.wip "!git add -A && git commit -m 'wip: save point'"
\end{shell}

A single word that snapshots everything before you try something risky. Phase 4
explains why this makes you recoverable, and Phase 5 explains how to squash the
\co{wip} commits away before anyone sees them.
\end{tryit}

\section{Ignoring files}

\subsection{Precedence}

Git checks several sources, most specific first:

\begin{enumerate}
  \item \co{.gitignore} in the same directory, then each parent directory
  \item \co{.git/info/exclude}, local to your clone and never committed
  \item The file named by \co{core.excludesFile}, your global ignore
\end{enumerate}

\subsection{Patterns}

\begin{terminal}
*.log              all .log files, any directory
build/             a directory anywhere
/build             build only at the repository root
!important.log     re-include something a previous rule excluded
doc/**/*.pdf       PDFs at any depth under doc/
temp?.txt          temp1.txt, tempA.txt
\end{terminal}

\begin{trap}
\co{.gitignore} only affects \textbf{untracked} files. Once a file is tracked, git
keeps tracking it no matter what you add to the ignore list. This is the single most
common confusion with the feature.

To stop tracking a file while keeping it on disk:

\begin{shell}
git rm --cached config.local.json
git commit -m "Stop tracking local config"
\end{shell}

The \co{--cached} is essential. Without it, \co{git rm} deletes the file.
\end{trap}

\subsection{A global ignore}

Editor and OS clutter belongs in your global ignore, not in every project's
\co{.gitignore}. Nobody else should have to accommodate your editor.

\begin{shell}
git config --global core.excludesFile ~/.gitignore_global
\end{shell}

Put \co{.DS\_Store}, \co{Thumbs.db}, \co{.idea/}, \co{*.swp} in there.

\section{gitattributes}

\co{.gitignore} decides which files git sees. \co{.gitattributes} decides how git
treats them. It is committed, so it applies to everyone.

\begin{terminal}
* text=auto eol=lf
*.png binary
*.pdf binary
*.md diff=markdown
*.py diff=python
package-lock.json linguist-generated=true
CHANGELOG.md merge=union
docs/ export-ignore
*.psd filter=lfs diff=lfs merge=lfs -text
\end{terminal}

Three of those are worth explaining.

\co{diff=python} tells git the language, so diff hunk headers show the enclosing
function name instead of a meaningless nearby line. Small change, and it makes
reviewing much easier.

\co{linguist-generated=true} collapses generated files in GitHub's diff view and
excludes them from the language statistics. Lock files are the usual case.

\co{merge=union} on a changelog makes git keep \emph{both} sides of a conflict
automatically instead of stopping. For an append-only list of lines this is usually
right, and it would have removed the conflicts from Phase 3, which is exactly why we
did not use it there.

\section{Hooks}

Scripts git runs at specific moments. They live in \co{.git/hooks} and there are
samples there already.

\begin{shell}
ls .git/hooks
\end{shell}

Remove the \co{.sample} suffix and make it executable to enable one.

\subsection{The useful client-side hooks}

\vspace{0.5em}
\begin{tabularx}{\linewidth}{@{}L{3.4cm} X@{}}
\toprule
\rhtablehead{Hook} & \rhtablehead{Runs} \\
\midrule
\co{pre-commit} & Before the message is written. Exit non-zero to cancel. Linting
and secret scanning go here. \\
\co{prepare-commit-msg} & Before the editor opens. Use it to insert a template or a
ticket number. \\
\co{commit-msg} & After the message is written. Enforce a format. \\
\co{pre-push} & Before a push. Run the test suite here, not in \co{pre-commit}. \\
\co{post-checkout} & After switching branches. Reinstall dependencies. \\
\bottomrule
\end{tabularx}
\vspace{0.5em}

A \co{pre-commit} hook that blocks obvious secrets:

\begin{terminal}
#!/bin/sh
if git diff --cached | grep -nE "(api[_-]?key|secret|password|BEGIN.*PRIVATE KEY)"; then
  echo "Possible secret in staged changes. Commit blocked."
  echo "If this is a false positive: git commit --no-verify"
  exit 1
fi
\end{terminal}

\begin{trap}
Hooks live in \co{.git/hooks}, which is not part of the repository and does not get
cloned. A hook you write protects you and nobody else.

Two ways round it. Point git at a directory that \emph{is} committed:

\begin{shell}
git config core.hooksPath .githooks
\end{shell}

Or use a framework that installs them on \co{npm install} or equivalent: Husky for
JavaScript, \co{pre-commit} for Python and general use.

Also remember \co{--no-verify} skips every hook. Hooks are a convenience that catches
mistakes, never a security control. Anything that genuinely must hold has to run in
CI, where the contributor cannot skip it. Phase 14 covers that.
\end{trap}

\subsection{Commit message conventions}

Conventional Commits is the widely used format:

\begin{terminal}
feat(parser): support blank lines in sonnets
fix: correct off-by-one in line counter
docs: explain the force-with-lease difference
refactor!: drop support for the old config format
\end{terminal}

The prefix is machine readable, which is the point: tools derive the next version
number and generate a changelog from it. \co{feat} bumps MINOR, \co{fix} bumps
PATCH, and \co{!} or a \co{BREAKING CHANGE:} footer bumps MAJOR.

Enforce it in \co{commit-msg} if your team wants it. Consider whether you want it
first, because on a small project it can be ceremony without benefit.

\section{Reference}

\vspace{0.5em}
\begin{tabularx}{\linewidth}{@{}L{6.6cm} X@{}}
\toprule
\rhtablehead{Goal} & \rhtablehead{Command} \\
\midrule
Where did this setting come from? & \co{git config -{}-list -{}-show-origin} \\
Different identity per directory & \co{includeIf "gitdir:\textasciitilde/work/"} \\
Push a new branch with no flags & \co{git config -{}-global push.autoSetupRemote true} \\
Make an alias & \co{git config -{}-global alias.lg "log -{}-graph ..."} \\
Alias that runs a shell command & prefix the value with \co{!} \\
Stop tracking a committed file & \co{git rm -{}-cached <file>} \\
Ignore editor clutter everywhere & \co{git config -{}-global core.excludesFile ...} \\
Share hooks with the team & \co{git config core.hooksPath .githooks} \\
Skip hooks once & \co{git commit -{}-no-verify} \\
\bottomrule
\end{tabularx}
\vspace{0.5em}

\checkyourself{
\begin{enumerate}
  \item Four config levels. Which wins?
  \item You added \co{config.json} to \co{.gitignore} and git still tracks it. Why?
  \item What does \co{.gitattributes} do that \co{.gitignore} cannot?
  \item Why can a \co{pre-commit} hook never be a security control?
  \item Which hook should run your test suite, and why not \co{pre-commit}?
\end{enumerate}}

\vspace{1em}
{\color{rhborder}\rule{\linewidth}{0.8pt}}

\textbf{Next:} Phase 12, \emph{GitHub as a platform}, which leaves git behind
entirely and covers the things GitHub added on top: issues, projects, pages, and a
search syntax worth learning properly.

\end{document}
